% 11-typecheck.tex — Chapter 11: Type Inference
\chapter{Type Inference}
\label{chap:11-typecheck}
\index{type inference}
\index{sema\_infer\_expr}

\pnum
Type inference is implemented in \srcref{src/sema/typecheck.h}.  It runs during
Phase~2 (\cname{walk\_stmt}) via \cname{sema\_infer\_expr(e)}, which visits an
expression and sets \cname{e->type}.

% -------------------------------------------------------------------------
\section{Builtin type singletons}
\label{sec:tc-builtins}
% -------------------------------------------------------------------------

\pnum
Two singleton constructors create canonical type nodes for built-in types:

\begin{itemize}
\item \cname{get\_builtin\_int\_type()} — returns a lazily-created \cname{Type*}
  for \lain{int} with \cname{base\_type->name = "int"}.
\item \cname{get\_builtin\_u8\_type()} — returns the canonical \lain{u8} type.
\end{itemize}

\pnum
These singletons are stored in \texttt{static} local variables.  Because all
compiler-internal type construction happens in \cname{sema\_arena}, the
singletons are created on first use and reused thereafter without re-allocation.

% -------------------------------------------------------------------------
\section{Integer widening}
\label{sec:tc-widening}
\index{integer widening}
% -------------------------------------------------------------------------

\pnum
Lain allows implicit widening of integer types in arithmetic expressions.
The widening rules are implemented by four helpers in
\srcref{src/sema/typecheck.h:58}:

\begin{center}
\begin{tabular}{ll}
\toprule
\textbf{Helper} & \textbf{Function} \\
\midrule
\cname{is\_integer\_type(t)}   & True for \lain{u8/i8/u16/i16/u32/i32/u64/i64/int/usize/isize} \\
\cname{integer\_rank(t)}       & Rank 1--5 (1=8-bit, 2=16-bit, 3=32-bit/int, 4=64-bit, 5=usize/isize) \\
\cname{wider\_integer\_type(a,b)} & Returns the type with the higher rank \\
\cname{can\_widen\_to(from,to)} & True if \texttt{from} rank $\leq$ \texttt{to} rank \\
\bottomrule
\end{tabular}
\end{center}

\pnum
For binary arithmetic and bitwise operators (\texttt{+}, \texttt{-}, \texttt{*},
\texttt{/}, \texttt{\%}, \texttt{\&}, \texttt{|}, \texttt{\^{}},
\texttt{<<}, \texttt{>>}), type inference returns the wider of the two
operand types.  Comparison and logical operators always return \lain{int}.
Narrowing (e.g.\ \lain{u64} to \lain{u8}) still requires an explicit
\lain{as} cast.  Float$\leftrightarrow$integer conversion also requires
explicit \lain{as}.

% -------------------------------------------------------------------------
\section{Type compatibility}
\label{sec:tc-compat}
\index{type compatibility}
% -------------------------------------------------------------------------

\pnum
\cname{types\_compatible(from, to)} (\srcref{src/sema/typecheck.h:98}) is a
conservative structural compatibility check used to diagnose argument/field type
mismatches.  It returns \texttt{true} for:
\begin{itemize}
\item \texttt{NULL} on either side (missing information; skip).
\item Same simple type name.
\item Integer widening (from rank $\leq$ to rank).
\item Pointer-to-same-element-type.
\end{itemize}
It returns \texttt{false} for clear mismatches (\lain{int} vs \lain{bool},
struct A vs struct B).

% -------------------------------------------------------------------------
\section{Inference for each ExprKind}
\label{sec:tc-exprkind}
% -------------------------------------------------------------------------

\pnum
\cname{sema\_infer\_expr(e)} dispatches on \cname{e->kind}:

\begin{center}
\begin{tabular}{ll}
\toprule
\textbf{Kind} & \textbf{Type assigned} \\
\midrule
\cname{EXPR\_LITERAL}        & \lain{int} \\
\cname{EXPR\_FLOAT\_LITERAL} & \lain{f64} \\
\cname{EXPR\_STRING}         & \lain{u8[:0]} (NUL-sentinel slice) \\
\cname{EXPR\_CHAR}           & \lain{u8} \\
\cname{EXPR\_BINARY}         & Wider operand type (arithmetic) or \lain{int} (comparison/logic) \\
\cname{EXPR\_UNARY}          & Operand type (or \lain{int} for logical not) \\
\cname{EXPR\_IDENTIFIER}     & Type from symbol table \\
\cname{EXPR\_CALL}           & Return type of callee declaration \\
\cname{EXPR\_MEMBER}         & Type of the named field in the struct declaration \\
\cname{EXPR\_INDEX}          & Element type of the array/slice \\
\cname{EXPR\_CAST}           & Target type of the cast \\
\cname{EXPR\_MOVE}           & Inner expression type with \cname{MODE\_OWNED} \\
\cname{EXPR\_MUT}            & Inner expression type with \cname{MODE\_MUTABLE} \\
\cname{EXPR\_ARRAY\_LITERAL} & Array type inferred from elements \\
\cname{EXPR\_RANGE}          & Range type (\lain{int..int}) \\
\bottomrule
\end{tabular}
\end{center}

% -------------------------------------------------------------------------
\section{String type inference and generics}
\label{sec:tc-strings}
% -------------------------------------------------------------------------

\pnum
String literals produce \lain{u8[:0]}, the NUL-sentinel slice type.  When a
string is passed to a comptime-generic function parameter that expects
\lain{comptime u8[:]}, the type is unified by the generic instantiation pass
(\cref{chap:18-comptime}) which creates a specialised \lain{Fixed\_u8\_N} type
for fixed-length strings.  The type inference layer emits \lain{Slice\_u8\_0}
for all runtime string values.

% -------------------------------------------------------------------------
\section{Parameter constraint propagation}
\label{sec:tc-constraints}
% -------------------------------------------------------------------------

\pnum
After parameter symbols are inserted, the orchestrator applies inline parameter
constraints (e.g.\ \lain{n int >= 0}) and pre-contract expressions to the
\cname{RangeTable}.  \cname{sema\_apply\_constraint(expr, ranges)} interprets a
binary comparison expression as a range restriction on the named variable and
calls \cname{range\_set} accordingly.  This seeds VRA with the proven parameter
bounds before any statement in the body is analysed.
